\section{The dimension-six specialization of the finite geometry}
\label{k6geom:section}

The finite optimization, named-endpoint geometry, and ordinary-suffix
construction also apply in dimension six. We verify this explicitly
rather than importing a dimension-specific formula.
Throughout this section $k=6$, $H=7$, $h_i=8-i$ for $1\leq i\leq6$,
and initially $L_i>0$. Original physical groups and original names
are retained, including when several consecutive products are equal.

\begin{proposition}[Finite geometric inputs at $k=6$]
\label{k6geom:finite}
The saddle, original transported coefficients, complete named endpoint
minima, and ordinary suffix in Section~\ref{profile:definition} are
well defined for $k=6$. They satisfy the same finite geometric estimates
used in the uniform proof, with constants depending only on $7$.
Neither an additional optimizer face nor a separate dimension-six
branch of the definition of $\Phi_k$ is required.
\end{proposition}

\begin{proof}
We check the ingredients and their boundary cases in turn.

\emph{The saddle and its compact range.}
The perspective recursion for $G_i$ is convex and increasing in its
child variable. Hence each $Z_i=e^{G_i}$ is convex, and $J$ is concave
on the closed ordered simplex. At a bottom plateau with product at
most $1/2$, raising that plateau has derivative bounded below by a
positive sum of terms $h-1-h^{1/2}\log h$. These are positive for
every $h>1$: with $t=\sqrt h$, the function
$t^2-1-2t\log t$ vanishes at $1$ and has derivative
$2(t-1-\log t)>0$. This also excludes zero products.

For a top plateau write $n$ for an exterior rank and $z$ for its
permanent-child weight. Its entropy increment is
\[
 E_n(z)=(n+z)\log(n+z)-z\log z,
 \qquad n\geq1,\quad z\geq1,\quad n+z\leq7.
\]
It satisfies $E_n(z)\geq2n\log2$. The derivative, with respect to
the exponent, of $(n+z)^s E_n(z)/(n+z)$ is at most
$7(\log7)^2$. Thus lowering a top product at least
$1-\epsilon_7$ improves $J$, where
\begin{equation}
 \epsilon_7=\frac{2\log2-1}{14(\log7)^2}>0.
 \label{k6geom:epsilon}
\end{equation}
Positivity here is analytic: $\log2=\int_1^2dt/t>1/2$.
Also $\epsilon_7<1/2$, for instance using $\log2<1$ and
$\log7>1$. We have proved $1/2<s_6\leq s_1<1-\epsilon_7$.

Uniqueness does not require six or more intermediate ranks. To see
this directly, set $S_6=2$ and $S_i=1+\exp(G_{i+1}/s_i)$ for $i<6$.
In a direction $v$ put $g_i=DG_i[v]$,
$Q_i=D^2G_i[v,v]$, and write $a_i=\log S_i$,
$u_i=s_{i+1}/s_i$, and $\beta_i=S_{i+1}^{u_i}/S_i$. Define
$A_i=a_i-\beta_i u_i a_{i+1}$ for $i<6$. Differentiation
gives
\[
 g_i=A_iv_i+\beta_i g_{i+1},\qquad
 Q_i=\beta_iQ_{i+1}
   +\frac{\beta_i(1-\beta_i)}{s_i}
                 (g_{i+1}-u_i a_{i+1}v_i)^2.
\]
At the terminal node $g_6=(\log2)v_6$ and $Q_6=0$.
All coefficients are compactly bounded on the product range just
proved. Since $g_i=a_iv_i+\beta_i(g_{i+1}-u_i a_{i+1}v_i)$,
backward induction yields
$D^2Z_i[v,v]\geq c_7\sum_{j\geq i}v_j^2$.
Consequently $-D^2J[v,v]\geq c_7\sum_j V_jv_j^2>0$ for
$v\ne0$. This proves uniqueness, also on every equality subspace.

\emph{Thresholds and genuine old flow.}
For every real $h>1$, set
$\alpha_h=\log((h-1)/\log h)/\log h$. With $x=(\log h)/2$,
\[
 \alpha_h=\frac12+\frac{\log(\sinh x/x)}{2x}.
\]
The second term is positive. Its derivative has the sign of
$N(x)=x\coth x-1-\log(\sinh x/x)$; here
$N(0+)=0$ and $N'(x)=1/x-x/\sinh^2x>0$.
Thus
\begin{equation}
 \frac12<\alpha_2<\alpha_3<\cdots<\alpha_7<1.
 \label{k6geom:alpha}
\end{equation}
The upper inequality follows from $\log h>1-1/h$.
Using the exact transported identity
\[
 Z_g\prod_{l=g}^{j-1}\beta_l
       =Z_j\prod_{l=g}^{j-1}\beta_l^{1-s_l}\leq Z_j
\]
in the bottom-plateau derivative strengthens the lower product bound
to $s_6>\alpha_2$. If that plateau contains a rank greater than two,
one native term is strictly positive; if it consists only of rank
two, the old transport inequality is strict and the old length is
positive. These exhaust the complete $7,6,\ldots,2$ chain.

Every adjacent product ratio exceeds $1/2$. A proper compound weight
therefore satisfies $z>\sqrt2$, whereas the terminal weight is exactly
$1$. Original old coefficients satisfy
\[
 7^{-8}\leq\kappa_g\leq1,\qquad
 \kappa_{g+1}-\kappa_g\geq d_*>0,
\quad
 d_*:=\min\left\{\frac{\epsilon_7}{49},
 7^{-8}\bigl[1-(6/7)^{\epsilon_7}\bigr]\right\}.
\]
These statements concern natural reference coefficients; they do not
resample an already revealed old list. Platform stationarity and the
entropy integral give the unchanged hierarchy
\[
 V_{a-1}\leq C_7\sum_{i=a}^bL_i,
 \qquad C_7=2\cdot7^3(1+\log7)^3.
\]
Indeed $E_p(x)/p\geq1+\log x+p/[2(x+p)]$ gives the incoming
surplus at least $np/[98(1+\log7)^2]$ at every strict platform
boundary. It is positive without a lower bound on the product gap.

\emph{KKT reference geometry and all named contexts.}
On any owner put $g_g(q)=q-\kappa_gE_q(z)$. The exact clipped
benchmarks satisfy
\[
 \Lambda_R=\sum_gL_gg_g(\min(r_g,R))
            =\lambda_{b-R}\geq0,
 \qquad \Lambda_0=\Lambda_{r_a}=0.
\]
For old groups the increasing coefficients permit their average
over the first region, with the nonnegative summation-by-parts
correction
$\int(\bar\kappa-\kappa(t))E_{d(t)}(z)\,dt$.
Reverse rank compression then gives each completed reference history
as a nonnegative sum of a clipped benchmark and compression losses.
The compression argument uses only $z\geq1$, integer ranks, and
the decreasing normalized secant
\[
 \frac{d}{dy}\frac{\log[(y\log y-b\log b)/(y-b)]}{\log y}
       \leq-\frac1{24\cdot49}\quad(1\leq b<y\leq7).
\]
For completeness, at $b=1$ this follows by writing
$F(t)=\log\int_0^1e^{tu}du$: the tilted density is at least $1/7$,
so $F''(t)\geq1/84$ and
$(tF'(t)-F(t))/(e^tt^2)\geq1/(24\cdot49)$.
Increasing $b$ decreases the displayed derivative, by the normalized-secant
lemma proved in Appendix~\ref{app:endpoint-details}.

All subsequent rank margins remain strictly positive. In particular
$E''(q)=1/(q+z)\geq1/7$, native $\kappa_g\geq1/7$, and
\[
 \Lambda_q+\frac{\log7}{2}\sum_{r_g=q}L_g
       \geq\frac1{14}\sum_{r_g>q}\kappa_gL_g
       \geq\frac{L_i}{98}\quad(r_i>q).
\]
This is the discrete second-difference identity with its two
nonnegative neighboring KKT margins; it does not assume global
concavity of $\Lambda$. The legal two-cut exchange pays a proper
input, the normalized-secant update and this same identity pay a
wrong retained suffix, and the true old average surplus pays an old
weak full input. Their thresholds can be chosen as the same finite
minima of positive dimension constants as in the general proof.
When a rank interval is empty its comparison is vacuous, rather than
requiring a nonexistent intermediate rank.

For the short-side endpoint estimate an explicit positive choice is
\[
 d_0=\min\left\{\frac{d_*}{\log7},\frac1{1176},
              \frac{\log(2\log2)}{7(\log7)^2}\right\},
 \qquad c_*=1-2^{-d_0/3}>0.
\]
For each clipped rank the normalized exponent increases by at least
$d_0$ between consecutive original groups. Encountering a coefficient
$f_g\leq c_*$ forces every later coefficient to be at most $-c_*$.
Together with $\Lambda_R\geq0$ this proves
$\min_R B_R(V_l)\geq c_*\min(V_l,V_b-V_l)$.
This assertion is only at original physical endpoints; a weak block
is not artificially split to obtain it.

The named enumeration in the profile consequently applies literally.
There are at most $21$ canonical edges and at most $7^6$ choices of
outside cuts for each edge. Each removed name has a deadline at least
the right endpoint of its assigned block. Even the left completion
has a cut at that block's left endpoint, so its last positive time is
at least $V_{a-1}$. This includes equality at the old/native boundary
and time zero for a top owner. The incoming-free score comparison is
therefore on its correct domain. It still subtracts the same incoming
constant from both endpoints and has the same exact binary increment;
all named contexts, not only deepest-suffix outside contexts, occur.

An owner of rank one causes no exception: a proper rank-one owner
has its one edge $1\to0$ with $z>\sqrt2$, and a terminal rank-one
owner has no genuine binary edge. Its zero-output edge is ordinary.
For every retained binary edge, if
$a=\log(r+z)$, $b=\log(q+z)$ and $p=(q+z)/(r+z)$, then
\[
 p,1-p\geq\frac17,\qquad b\geq\frac{\log2}{2},
 \qquad \frac17\leq a-bp\leq\log7.
\]
Thus its two-mark variance is positive uniformly, as is its mean,
and its natural $\rho$ lies in a fixed compact subinterval of $(0,1)$.
No degenerate terminal $q=0$ law is inserted in the binary product.

The original branching comparisons also retain positive margins.
For a proper owner one can take
$\gamma=(\sqrt2-1)\log(8/7)>0$ in the pure-cell entropy gain,
because its compound mass is at least $\sqrt2$ and at most $7$.
If the first name of that cell has expired, the gain is instead at
least $2\log2$. For the terminal owner, where the proper-child gap
cannot be used, a floating nonsingleton chain has
$1\leq d(t)\leq r(t)-1$ throughout its floating interval. Its two
legal reference rank changes $d\pm1$ give the lower bound
$7^{-8}T_{\rm float}/14$ for its reference cost. At $r=1$ such an
interval is empty. Mixed-time domination uses only eligible-rank
addition under the lawful parent/child projections. Thus the same
positive mixed, spread, and reference shares pay the original birth
errors and erased-pivot fibers with at most six original boundaries;
no extra small-rank case is omitted. These are deterministic comparisons
under the same named regularity and allocated-wall hypotheses as in
the general history theorem.

\emph{The ordinary suffix.}
Set $c_w=1/196$ and use exactly
\[
 \tau_7=\min\left\{\frac14,
    \frac1{2744(\log7)^2},
    \frac14\min_{2\leq h<7}(\alpha_{h+1}-\alpha_h)\right\}>0.
\]
The strict positivity of every gap was proved in
\eqref{k6geom:alpha}; no numerical gap test is used.
For $e_h=h^\sigma\log h-(h-1)$,
$|e_h'|\leq7(\log7)^2$, so $|e_h|\leq c_w/2$ in its
selected threshold interval. The intervals are disjoint.
Moreover a genuine terminal binary weak transition $r\to q\geq1$
at threshold $\epsilon_w\leq1/(4\cdot7^3)$ obeys
$g(r)\geq1/196$. Indeed strong concavity gives
\[
 g(q)-\frac qr g(r)\geq\frac{q(r-q)}{98},
 \qquad g(r)\geq\frac{qr}{98}-7\epsilon_w\geq\frac1{196}.
\]
Its ordinary drift is therefore at most $-c_w$, so it cannot occur
in the selected ordinary suffix.

The actual terminal platform balance gives total ordinary mean
$D_{\rm all}\geq b_7V_{\rm old}$, where
$b_7=1/[98(1+\log7)^2]$; for empty old input it gives zero.
All groups before the selected suffix have drift at most a fixed
$-g_0<0$, because their finitely many thresholds are separated by
$\tau_7$. Hence
\[
 D_{\rm suffix}\geq b_7V_{\rm old}+g_0L_{\rm head}\geq0,
 \qquad U\geq c_7V_6.
\]
A positive-drift block exists whenever no present block is in a
threshold interval, since otherwise the total mean would be negative.
There are at most six suffix blocks, so its largest has length at
least $U/6$. The empty-prefix, single-block, and zero total mean cases
all remain included. This is precisely the geometric input to the
single ordinary root, with the original counts and gaps unchanged.
\end{proof}

\begin{proposition}[Explicit lift envelope at dimension six]
\label{k6geom:lift}
The finite search and actual-parameter-stability comparison in the
definition of $\Phi_k$ have positive constants at $k=6$. Their values are
\[
 A_6=686,\quad B_6=174,\quad N=C=21,\quad E=25,
 \quad c=42,\quad a=A_6-C=665.
\]
In particular a valid search bound is
\[
 R=1+\max\left\{1,\max_iL_i,
       \frac{105000}{442225},
       \frac{2(687+5\log2)}{665}\right\}.
\]
The same formula for $\Phi_k$ therefore includes $k=6$ on all
$\ell_i\geq\log3$, including zero original block lengths.
\end{proposition}
\begin{proof}
Substitution in the displayed dimension formulas gives these exact
values. In particular $B_6<A_6$ and $a>0$. There are at most $21$
binary factors, each bounded below by $(1+7L_i)^{-1}$;
there are at most five proper exit factors, with $0<d_j<1$; and
the sole ordinary factor is at least $(1+V_6)^{-3/2}$.
For $\lambda\geq\max_iL_i$ the lifted objective is thus at least
$2^{-5}e^{665\lambda}(1+42\lambda)^{-25}$.
The inequality $\log(1+x)\leq\sqrt x$ gives the stated finite
range exactly as in the profile proof. The actual arithmetic
stability exponent is $174\lambda+O_6(1)$, strictly less than the
chosen envelope penalty. Positive lifted lengths suffice for the
geometry above; the infimum and actual stability, rather than an
assumption about optimizer-face continuity, recover zero blocks.
\end{proof}

\begin{remark}[Scope of the specialization]
This verification is of the finite geometric and defining inputs.
The scalar bridge, full-type conditioning, original-class row, and
arithmetic transfer steps keep their stated hypotheses and their
separate proofs. None acquires independence by setting $H=7$.
Once those dimension-compact inputs are invoked, the same upper,
lower, and lift-removal proof gives the dimension-six case using the
same $\Phi_k$; no use of a separate older $k=6$ profile is needed.
\end{remark}
